\begin{abstract}
As multi-turn conversations grow longer, large language models progressively lose attention to the system-prompt goals that define their task, even when the underlying goal representations remain intact. Existing remedies operate at inference time, leaving a natural but untested alternative: training auxiliary attention losses alongside next-token prediction to directly optimize goal attention. This strategy assumes that attention patterns learned under teacher forcing will transfer to autoregressive generation, an assumption that classical exposure bias research has challenged for sequence models but that has never been examined for auxiliary attention objectives in multi-turn transformers.
We show this assumption fails. On a controlled five-way testbed spanning two architectures (Qwen-2.5-7B, Llama-3.1-8B), goal attention optimized to target under teacher forcing collapses by an order of magnitude under autoregressive generation. A diagnostic exclusion chain traces the failure to exposure bias, ruling out gradient conflict, gating mechanism, and parameter capacity as primary causes. Inference-time attention steering confirms that the collapse reflects attention misallocation rather than representation loss. Scheduled sampling, applied here for the first time to auxiliary attention objectives in multi-turn transformers, compresses the training--inference gap to near parity on Qwen and from ${\sim}$27$\times$ to 2.9$\times$ on Llama, with no perplexity cost\footnote{Code and data are available at \url{https://anonymous.4open.science/r/turn_residual_gating-8703/}.}.
\end{abstract}
